% ============================================================================
% LANGENTWURF-TEMPLATE - Digitale Grundlagen (Senioren 65+)
% ============================================================================
% Basiert auf: Greenhouse Schools Langentwurf-Template v2.0
% Angepasst für: Erwachsenenbildung Senioren 65+
% Profil: S (Senioren)
% Tier: 2 (Standard)
%
% VERWENDUNG:
%   1. Kopiere diese Datei als LE-XXX-thema.tex
%   2. Fülle die Variablen aus
%   3. Kompiliere 2x mit pdflatex
% ============================================================================

\documentclass[11pt,a4paper]{article}

% ============================================================================
% PACKAGES
% ============================================================================
\usepackage[utf8]{inputenc}
\usepackage[T1]{fontenc}
\usepackage[ngerman]{babel}
\usepackage{geometry}
\usepackage{fancyhdr}
\usepackage{lastpage}
\usepackage{xcolor}
\usepackage{tcolorbox}
\usepackage{enumitem}
\usepackage{tabularx}
\usepackage{longtable}
\usepackage{array}
\usepackage{booktabs}
\usepackage{hyperref}
\usepackage{tikz}
\usepackage{amssymb}

% ============================================================================
% KONFIGURATION (FEST FÜR PROFIL S)
% ============================================================================
\newcommand{\plantier}{2}           % Tier 2 = Standard
\newcommand{\profil}{S}             % S = Senioren 65+
\newcommand{\fachfarbe}{b35c00}     % Orange (Informatik/Digital)

% ============================================================================
% VARIABLEN - HIER AUSFÜLLEN
% ============================================================================

% --- Metadaten ---
\newcommand{\einheitnummer}{001}                    % 001-023
\newcommand{\einheitthema}{iPhone-Grundlagen}       % Konkretes Thema
\newcommand{\modulname}{Geräte \& Zugänge}          % A, B, C, D oder E
\newcommand{\modulbuchstabe}{A}

% --- Lernziele (3-5 Hauptziele) ---
\newcommand{\lernzielA}{Die drei Zugangscodes (SIM-PIN, Geräte-Code, Apple-ID) unterscheiden}
\newcommand{\lernzielB}{Die Schutzfunktion jedes Codes verstehen}
\newcommand{\lernzielC}{Wissen, was bei Verlust eines Codes passiert}

% --- Zeitrahmen ---
\newcommand{\zeitrahmen}{60--90 Minuten}            % Flexible Kurseinheit

% --- Autor ---
\newcommand{\autor}{N.N.}

% ============================================================================
% LAYOUT
% ============================================================================
\geometry{left=2.5cm, right=2.5cm, top=2.5cm, bottom=2.5cm}
\definecolor{fach}{HTML}{\fachfarbe}
\definecolor{gray}{HTML}{666666}
\definecolor{lightgray}{HTML}{f5f5f5}
\definecolor{successgreen}{HTML}{2a7a4b}
\definecolor{warningorange}{HTML}{b35c00}
\definecolor{errorred}{HTML}{c62828}

\tcbuselibrary{skins,breakable}

% Farbige Boxen
\newtcolorbox{infobox}[1][]{
    colback=lightgray,
    colframe=fach,
    fonttitle=\bfseries,
    title=#1,
    breakable
}

\newtcolorbox{warnbox}[1][]{
    colback=errorred!5,
    colframe=errorred,
    fonttitle=\bfseries,
    title=#1,
    breakable
}

\newtcolorbox{scorebox}{
    colback=white,
    colframe=fach,
    fonttitle=\bfseries,
    title=Didaktik-Score,
    breakable
}

% Kopf-/Fußzeile
\pagestyle{fancy}
\fancyhf{}
\fancyhead[L]{\small\textcolor{gray}{Digitale Grundlagen -- \einheitnummer{} \einheitthema}}
\fancyhead[R]{\small\textcolor{gray}{Langentwurf Tier \plantier{} / Profil \profil}}
\fancyfoot[C]{\small\thepage{} / \pageref{LastPage}}
\renewcommand{\headrulewidth}{0.4pt}
\renewcommand{\footrulewidth}{0pt}

% Abschnitte
\renewcommand{\thesection}{\arabic{section}}
\renewcommand{\thesubsection}{\thesection.\arabic{subsection}}

% ============================================================================
% DOKUMENT
% ============================================================================
\begin{document}

% ============================================================================
% DECKBLATT
% ============================================================================
\begin{titlepage}
    \centering
    \vspace*{2cm}

    {\Large\textcolor{gray}{Digitale Grundlagen}}\\[0.5cm]
    {\large\textcolor{gray}{Erwachsenenbildung für Senioren 65+}}\\[2cm]

    {\Huge\textcolor{fach}{\textbf{Langentwurf}}}\\[0.5cm]
    {\large Tier 2 (Standard) \quad|\quad Profil S (Senioren)}\\[2cm]

    {\LARGE\textbf{\einheitnummer{} -- \einheitthema}}\\[0.5cm]
    {\large Modul \modulbuchstabe: \modulname}\\[2cm]

    \begin{tabular}{rl}
        \textbf{Zeitrahmen:} & \zeitrahmen \\[0.3cm]
        \textbf{Zielgruppe:} & Senioren 65+ \\[0.3cm]
        \textbf{Vorkenntnisse:} & Minimal bis keine \\[0.3cm]
    \end{tabular}

    \vfill

    {\large Autor: \autor}\\[0.3cm]
    {\normalsize\textcolor{gray}{Stand: \today}}\\[1cm]

\end{titlepage}

% ============================================================================
% INHALTSVERZEICHNIS
% ============================================================================
\tableofcontents
\newpage

% ============================================================================
% KONFIGURATIONSÜBERSICHT
% ============================================================================
\begin{infobox}[Konfiguration]
\begin{tabular}{ll}
    \textbf{Einheit:} & \einheitnummer{} -- \einheitthema \\[0.3em]
    \textbf{Modul:} & \modulbuchstabe{} -- \modulname \\[0.3em]
    \textbf{Tier:} & 2 (Standard) \\[0.3em]
    \textbf{Profil:} & S (Senioren 65+) \\[0.3em]
    \textbf{Zeitrahmen:} & \zeitrahmen \\[0.3em]
    \textbf{Differenzierung:} & Basis / Vertiefung (keine Noten)
\end{tabular}
\end{infobox}

% ============================================================================
% PROFIL S: SENIOREN 65+
% ============================================================================
\begin{infobox}[Profil S: Senioren 65+]
\begin{tabular}{ll}
    Zielgruppe: & Erwachsene 65+, ``digitale Analphabeten'' \\
    Vorkenntnisse: & Minimal bis keine \\
    Lernziel: & Überleben in der digitalen Welt \\
    Prüfungen: & Keine (Selbst-Checklisten) \\
    Besonderheiten: & Große Schrift, Alltagsanalogien, Sicherheitsfokus
\end{tabular}

\vspace{0.5em}
\textbf{Didaktische Prinzipien:}
\begin{itemize}[nosep]
    \item Alltagsanalogien statt Technik-Jargon
    \item ``Überleben'' statt ``Meistern''
    \item Sicherheitsaspekte betonen
    \item Vertrauenspersonen einbeziehen
\end{itemize}
\end{infobox}

% ============================================================================
% 1. BEDINGUNGSANALYSE
% ============================================================================
\section{Bedingungsanalyse}

\subsection{Zielgruppe}
\begin{itemize}
    \item \textbf{Alter:} 65+ Jahre
    \item \textbf{Digitale Vorerfahrung:} Minimal bis keine
    \item \textbf{Motivation:} Praktische Notwendigkeit (Kommunikation mit Familie, Behörden, Banking)
    \item \textbf{Ängste:} Oft vorhanden (``Ich mache etwas kaputt'', ``Das ist zu kompliziert'')
\end{itemize}

\subsection{Besondere Anforderungen}
\begin{itemize}
    \item \textbf{Sehvermögen:} Große Schrift ($\geq$11pt), hoher Kontrast
    \item \textbf{Motorik:} Große Klickflächen, langsames Tempo
    \item \textbf{Gedächtnis:} Wiederholung, Checklisten zum Nachschlagen
    \item \textbf{Sprache:} Einfach, keine Anglizismen ohne Erklärung
\end{itemize}

\subsection{Lernumgebung}
\begin{itemize}
    \item \textbf{Setting:} Einzelunterricht oder Kleingruppe (max. 5 Personen)
    \item \textbf{Geräte:} Eigene Geräte der Teilnehmer (iPhone, iPad, Mac)
    \item \textbf{Materialien:} Gedruckte Unterlagen zum Mitnehmen
    \item \textbf{Begleitung:} Ggf. Vertrauensperson anwesend
\end{itemize}

% ============================================================================
% 2. SACHANALYSE
% ============================================================================
\section{Sachanalyse}

\subsection{Kernkonzepte der Einheit}

% TODO: Für jede Einheit anpassen
\begin{enumerate}
    \item \textbf{Konzept 1:} ...
    \item \textbf{Konzept 2:} ...
    \item \textbf{Konzept 3:} ...
\end{enumerate}

\subsection{Alltagsanalogien}

% TODO: Für jede Einheit anpassen
\begin{tabular}{|l|l|}
\hline
\textbf{Fachbegriff} & \textbf{Alltagsanalogie} \\
\hline
... & ... \\
\hline
\end{tabular}

\subsection{Typische Missverständnisse}

% TODO: Für jede Einheit anpassen
\begin{itemize}
    \item ...
\end{itemize}

% ============================================================================
% 3. DIDAKTISCHE ANALYSE
% ============================================================================
\section{Didaktische Analyse}

\subsection{Stellung in der Gesamtreihe}

Dies ist Einheit \textbf{\einheitnummer} im Modul \textbf{\modulbuchstabe: \modulname}.

% TODO: Einordnung in Gesamtkonzept

\subsection{Legitimation}

\textbf{Alltagsrelevanz:}
\begin{itemize}
    \item Warum muss ein Senior das wissen?
    \item Was passiert, wenn er es nicht weiß?
\end{itemize}

\textbf{Sicherheitsaspekte:}
\begin{itemize}
    \item Welche Risiken werden adressiert?
    \item Welche Schutzmaßnahmen werden vermittelt?
\end{itemize}

\subsection{Didaktische Reduktion}

\textbf{Bewusst ausgeklammert:}
\begin{itemize}
    \item ... (zu komplex für Einsteiger)
    \item ... (nicht überlebensnotwendig)
\end{itemize}

\textbf{Fokus auf:}
\begin{itemize}
    \item Das Minimum zum ``Überleben''
    \item Sicherheitskritische Aspekte
\end{itemize}

% ============================================================================
% 4. LERNZIELE
% ============================================================================
\section{Lernziele}

\subsection{Grobziel}

Die Teilnehmer verstehen \textbf{[Thema]} und können \textbf{[Handlung]} selbstständig durchführen.

\subsection{Feinziele}

\begin{tabular}{|c|p{8cm}|c|c|}
\hline
\textbf{Nr.} & \textbf{Die Teilnehmer können...} & \textbf{Stufe} & \textbf{Niveau} \\
\hline
FZ1 & \lernzielA & Wissen & Basis \\
\hline
FZ2 & \lernzielB & Verstehen & Basis \\
\hline
FZ3 & \lernzielC & Anwenden & Basis \\
\hline
FZ4 & ... & Anwenden & Vertiefung \\
\hline
\end{tabular}

\vspace{1em}
\textbf{Taxonomie-Verteilung (angepasst für Senioren):}
\begin{itemize}
    \item Wissen (erinnern, nennen): ca. 40\%
    \item Verstehen (erklären, zuordnen): ca. 40\%
    \item Anwenden (durchführen, lösen): ca. 20\%
\end{itemize}

% ============================================================================
% 5. METHODISCHE ANALYSE
% ============================================================================
\section{Methodische Analyse}

\subsection{Methodische Grundsätze}

\begin{itemize}
    \item \textbf{Tempo:} Langsam, mit Pausen zur Verarbeitung
    \item \textbf{Wiederholung:} Kernpunkte mehrfach in verschiedenen Formen
    \item \textbf{Aktivierung:} Sofort am eigenen Gerät ausprobieren
    \item \textbf{Sicherheit:} ``Das können Sie nicht kaputt machen''
\end{itemize}

\subsection{Sozialformen}

\begin{tabular}{|l|l|l|}
\hline
\textbf{Phase} & \textbf{Sozialform} & \textbf{Begründung} \\
\hline
Einführung & Plenum/Einzeln & Gemeinsamer Fokus \\
Erklärung & Demonstration & Zeigen am Gerät \\
Übung & Einzeln mit Begleitung & Individuelles Tempo \\
Sicherung & Gemeinsam & Checkliste durchgehen \\
\hline
\end{tabular}

\subsection{Materialien}

\begin{enumerate}
    \item \textbf{Infoblatt (IB):} 1-2 Seiten, große Schrift, Kernwissen
    \item \textbf{Übungen (UE):} Zuordnung, Richtig/Falsch, Situationen
    \item \textbf{Checkliste (CL):} ``Das habe ich verstanden''
    \item \textbf{Notfall-Notizen (NN):} Zum Ausfüllen und Aufbewahren
\end{enumerate}

\subsection{Differenzierung}

\begin{tabular}{|l|p{5cm}|p{5cm}|}
\hline
& \textbf{Basis} & \textbf{Vertiefung} \\
\hline
\textbf{Inhalt} & Überlebensminimum & Zusätzliche Optionen \\
\hline
\textbf{Übungen} & Zuordnung, Ja/Nein & Situationen, Transfer \\
\hline
\textbf{Ziel} & ``Ich schaffe das Wichtigste'' & ``Ich verstehe auch warum'' \\
\hline
\end{tabular}

% ============================================================================
% 6. VERLAUFSPLANUNG
% ============================================================================
\section{Verlaufsplanung}

\begin{longtable}{|p{1.2cm}|p{2cm}|p{4.5cm}|p{3.5cm}|p{2cm}|}
\hline
\textbf{Zeit} & \textbf{Phase} & \textbf{Inhalt} & \textbf{TN-Aktivität} & \textbf{Material} \\
\hline
\endhead

0--10 & Einstieg & Begrüßung, Thema vorstellen, Alltagsbezug herstellen & Zuhören, Fragen stellen & -- \\
\hline

10--25 & Erklärung & Kernkonzepte mit Analogien erklären, am Gerät zeigen & Zuhören, Beobachten & Infoblatt \\
\hline

25--45 & Übung & Am eigenen Gerät ausprobieren, Übungen bearbeiten & Selbst durchführen & Übungen \\
\hline

45--55 & Sicherung & Checkliste durchgehen, Fragen klären & Abhaken, Nachfragen & Checkliste \\
\hline

55--60 & Abschluss & Notfall-Notizen ausfüllen, Ausblick & Notieren & Notfall-Notizen \\
\hline
\end{longtable}

\textit{Hinweis: Zeiten sind flexibel. Bei Bedarf Pausen einlegen.}

% ============================================================================
% 7. ANLAGEN
% ============================================================================
\section{Anlagen}

\begin{enumerate}
    \item Infoblatt (IB-\einheitnummer-...)
    \item Übungen (UE-\einheitnummer-...)
    \item Checkliste (CL-\einheitnummer-...)
    \item Notfall-Notizen (NN-\einheitnummer-...)
\end{enumerate}

% ============================================================================
% 8. LÖSUNGEN
% ============================================================================
\section{Lösungen}

\subsection{Lösungen Übungen}

% TODO: Für jede Einheit ausfüllen

\subsection{Erwartete Fragen}

\begin{itemize}
    \item \textbf{Frage:} ...
    \item \textbf{Antwort:} ...
\end{itemize}

% ============================================================================
% 9. LITERATUR
% ============================================================================
\section{Literatur}

\begin{itemize}
    \item Apple Support: \url{https://support.apple.com/de-de}
    \item ...
\end{itemize}

% ============================================================================
% 10. DIDAKTIK-SCORE
% ============================================================================
\newpage
\section{Didaktik-Score}

\begin{scorebox}
\textbf{Bewertung der didaktischen Qualität -- Profil S (Senioren)}\\
\textbf{Ziel-Score: $\geq$ 9.0}

\vspace{1em}
\begin{tabular}{|c|l|c|c|p{4cm}|}
\hline
\textbf{\#} & \textbf{Dimension} & \textbf{Gew.} & \textbf{Score} & \textbf{Notiz} \\
\hline
1 & Differenzierung (Basis/Vertiefung) & 15\% & /10 & \\
\hline
2 & Taxonomie (Wissen/Verstehen/Anwenden) & 10\% & /10 & \\
\hline
3 & Alltagsrelevanz & 10\% & /10 & \\
\hline
4 & Aktivierung (Hands-on) & 10\% & /10 & \\
\hline
5 & Scaffolding (Schritt-für-Schritt) & 10\% & /10 & \\
\hline
6 & Feedback \& Ermutigung & 10\% & /10 & \\
\hline
7 & Sprachliche Klarheit (keine Jargon) & 10\% & /10 & \\
\hline
8 & Motivation (``Warum muss ich das?'') & 10\% & /10 & \\
\hline
9 & Barrierefreiheit (Schrift, Kontrast) & 10\% & /10 & \\
\hline
10 & Praktikabilität (Zeit, Material) & 5\% & /10 & \\
\hline
\multicolumn{3}{|r|}{\textbf{GESAMT:}} & \textbf{/10} & \\
\hline
\end{tabular}

\vspace{1em}
\textbf{Bewertungsschwellen:}
\begin{itemize}
    \item[$\Box$] \textcolor{successgreen}{$\geq$ 9.0: \textbf{FREIGABE}}
    \item[$\Box$] \textcolor{warningorange}{8.0--8.9: Iteration empfohlen}
    \item[$\Box$] 6.0--7.9: Überarbeitung erforderlich
    \item[$\Box$] $<$ 6.0: Grundlegende Neukonzeption
\end{itemize}
\end{scorebox}

\subsection{Verbesserungsmaßnahmen}

Falls Score $<$ 9.0:

\begin{enumerate}
    \item \textbf{[Schwächste Dimension]:} ...
    \item \textbf{[Zweitschlechteste]:} ...
    \item \textbf{[Drittschlechteste]:} ...
\end{enumerate}

% ============================================================================
% ENDE
% ============================================================================
\end{document}
